
Desde el origen de los primeros sistemas operativos, ha habido tareas que han demandado una ejecución precisa respecto a ciertas restricciones temporales. De esta necesidad, nacen los sistemas operativos de tiempo real (RTOS), donde se pueden implementar herramientas especializadas para gestionar dichas restricciones. Los circuitos biohíbridos formados por neuronas vivas y modelos neuronales que imitan su comportamiento son un ejemplo de sistema que ha de implementarse con software o hardware de tiempo real \cite{Amaducci2019,Amaducci2022}.

Los circuitos biohíbridos tienen múltiples aplicaciones en neurociencia y en neurotecnología, pues brindan nuevas posibilidades para el estudio de la dinámica neuronal, y también abren nuevas perspectivas en el campo sanitario y de neurorrehabilitación permitiendo expandir el conocimiento y la comprensión del funcionamiento del sistema nervioso. 

No obstante, la conectividad neuronal biohíbrida presenta diversos retos, entre ellos el intercambio de información compatible y eficaz entre ambas partes. La diferencia entre los dos medios en los que la entidad viva y la entidad artificial residen, al ser uno biológico y otro digital (o electrónico), es un factor clave, pero también lo es la diferencia temporal en la que operan estas entidades. Los procesadores no siguen de forma continua la cronología de una neurona viva, pues mientras que en esta los tiempos no son estrictamente constantes (debido a distintos factores que pueden influenciar a la neurona), en los procesadores se sigue un \textit{clock} o reloj, que sincroniza sus acciones. Añadido a esto, está el planificador del sistema operativo, que a través de las interrupciones, puede afectar al flujo de respuesta del modelo artificial, tanto si este está siendo ejecutado en el propio ordenador  \cite{Amaducci2019, Cerqueira2013}, como si es tan solo un medio para conectar un \textit{modelo neuronal} (como el de Hodgkin–Huxley \cite{Hodgkin1990}) a la neurona viva. Para paliar con esto último, se usan los Sistemas Operativos de Tiempo Real (\textit{RTOS}).

En los \textit{RTOS}, se pueden generar procesos que no respeten el curso natural del planificador, es decir, que tengan más precisión temporal de la que tendrían si fuese un Sistema Operativo de Propósito General (\textit{GPOS}). Esto permite que dicho proceso realice tareas con un margen temporal mucho más estrecho que un \textit{GPOS}, pues puede acabar de realizar la tarea antes de ser interrumpido, o dicho de otra forma, no se le interrumpe hasta que no acabe la tarea o hasta que otro proceso de mayor prioridad sea ejecutado. Por ello, estos sistemas son una buena opción para resolver el problema anteriormente planteado. 

Hay diversos factores a tratar sobre el sistema a elegir, a pesar de que se repasarán los más usados actualmente, el enfoque principal es el parche de Linux \textit{Preempt-RT} (concretamente en Debian), debido a su reciente fusión en \textit{mainline}. Este evento reciente, no solo marca un estándar en estos tipos de sistemas, sino que también facilita la instalación de un \textit{RTOS}, abriendo la puerta a que nuevos usuarios puedan desarrollar nuevas herramientas que faciliten o extiendan su uso, particularmente en el contexto de la neurotecnología.

